% Status: drafted 2026-09-22.
%
% Every count in this section is a macro rendered from eval/results/corpus.csv,
% training.csv and throughput.csv. Three figures the outline planned for this
% section are absent on purpose: the baseline text's chapter and word-token
% counts, and the size of the seed index. They are properties of a production
% database this repository does not hold, so PROVENANCE.md lists them among the
% figures that are not reproduced here rather than the report stating them.
\section{The corpus}
\label{sec:corpus}

This report measures one pipeline over two bodies of sixteenth-century French
print that are not independent of one another. The first is
\emph{Amadis de Gaule}, a chivalric romance cycle of \ocrBooks{} books running
to \ocrPages{} scanned pages. The second is the \emph{Trésor} anthologies,
printed compilations that quarry extracts out of that cycle and reprint them
under thematic headings. An extract in a \emph{Trésor} is therefore a passage
of \emph{Amadis} that already exists elsewhere in the same corpus, which is
what turns the editorial question of \emph{where does this piece come from}
into a retrieval problem with a checkable answer. The design specification
settles both works as one target domain rather than as a domain and an
out-of-domain control, and this report follows that: the two are reported
separately where the type differs, and never as in-domain against
out-of-domain.

\subsection{Two editions, two type families}

The \ocrBooks{} books do not form one edition. The first \ocrBooksFamilyA{}
books, called family~A throughout this report, are a large folio setting in
\emph{bâtarde} type with roman folio numbers and a printed table of contents at
the front of each volume. The remaining \ocrBooksFamilyB{} books, family~B,
are a later and smaller edition in clean roman type, with italic rubrics,
arabic page numbers, and no printed table of contents at all. The two families are of very
different weight: family~A is \ocrPagesFamilyA{} pages against family~B's
\ocrPagesFamilyB{}, so two thirds of the corpus is the later setting.

The division is not a presentational detail, and three later decisions rest on
it. Recognition is reported per family, because a model fine-tuned on one kind
of type has no claim on the other until it is shown. Structural calibration
exists for family~A alone, because calibration needs a printed table of
contents to score against and family~B prints none; section~\ref{sec:results}
reports no corpus-wide structural figure for that reason. Throughput, by
contrast, pools the two, because a page is a page to the recogniser and every
book was run by one model on one device at one working width.

\subsection{Scan provenance}

The page images reach the pipeline from three holding institutions and in
three different conditions. Book~1 arrives as an archive of bitonal 600~dpi
TIFFs from the Bibliothèque municipale de Lyon. Books~3 and~4 are greyscale
PDFs from Gallica. Book~18 is a copy held by the Biblioteca Valenciana and
carries that library's stamp on the page, which the layout stage suppresses as
an illustration region rather than reading it as text. The provenance of the
remaining volumes is not recorded in any artefact this repository holds, and
is therefore not claimed here.

Provenance matters to the evaluation twice over. It is a source of variation
that has nothing to do with the century the text was printed in, so
section~\ref{sec:protocol} stratifies the gold set on it and the stamped copy
is its own stratum. It also constrains what this repository may contain: reuse
terms differ per holding institution and are unsettled for at least one of the
three, so no page image and no transcribed text is committed here. What is
committed is counts, timings and offsets, which is enough to rerun every
evaluation in this report and not enough to reconstitute the corpus.

\subsection{The \emph{Trésor} pieces and the localisation reference}

A \emph{Trésor} piece is a contiguous span of the baseline text. The reference
for where each piece belongs is the editor's own catalogue, transcribed into
\file{data/localisation/ground-truth.csv}: \corpusPieces{} pieces, each with a
Livre and, where the print gives one, a chapter. The catalogue is prose written
by a scholar rather than a clean database field, so the transcription keeps the
editor's own hedging instead of flattening it. \corpusCertain{} pieces
(\corpusCertainShare{}) carry an unqualified assignment; \corpusUncertain{}
are marked with a question mark, \corpusConjecture{} are bracketed as
conjectures, and \corpusUnassigned{} name no location at all. Only the
\corpusCertain{} certain pieces are scored in section~\ref{sec:results}, and
the other three bands are reported as what they are: material the editor did
not commit to, which a matcher cannot be graded against.

The pieces are drawn from \corpusLivres{} of the \ocrBooks{} books, so the
anthologies are selective rather than a systematic abridgement of the cycle.
They are also short against what they are searched for. The matcher placed
\corpusAligned{} of the \corpusPieces{} pieces and recorded the span of each:
the median piece is \corpusCharsMedian{} code points, the mean
\corpusCharsMean{}, the shortest \corpusCharsMin{} and the longest
\corpusCharsMax{}. The search space each of them is sought in is the full
transcription of all \ocrBooks{} books. Its size in words is a property of the
production database, and no artefact in this repository records it, so this
report does not state one.

The pieces arrive in two batches, and the report keeps them apart everywhere.
The workbook cohort is the \corpusPiecesWorkbook{} pieces taken from the
editor's workbook, which is the batch that existed when E5's decision rule was
pre-registered. The catalogue cohort is the \corpusPiecesCatalogue{} pieces
taken later from a second document, the \emph{Juxtalinéaire} catalogue, and
imported four days afterwards. Pooling them would mean
enlarging a test set after seeing the numbers it produced, so they are never
pooled, and each carries its own denominators. The two batches also differ in
the material itself and not only in when they arrived: the workbook's pieces
average \corpusCharsMeanWorkbook{} code points against the catalogue's
\corpusCharsMeanCatalogue{}, about a third longer, which is one reason a single
figure over the two would be hard to read even if it were legitimate to compute.

\subsection{The training corpus}

The fine-tune saw \trainPages{} annotated pages, and every one of them comes
from the same Transkribus collection, \emph{Trésor des Amadis}~T.1. Of those,
\trainPagesTrain{} were used for training and \trainPagesVal{} were held
out. The held-out \trainPagesVal{} are validation and not test: they selected
the shipped checkpoint, and section~\ref{sec:limitations} states what that
costs. The split is stated here as a re-derived fact rather than an inherited
one. The original file lists no longer exist, so the split was recomputed from
the two export archives and confirmed against the compiled training data;
\file{docs/notes/2026-09-20-artefact-recovery.md} records the derivation and
the control that rules out a different random seed.

Two consequences follow, and both cut against the system. The first is that
the training material is one anthology volume, so both type families of
\emph{Amadis de Gaule} itself are largely unseen by the model that transcribed
them, and whether the fine-tune carries across the families is an open question
this report answers rather than assumes. The second is subtler. No page of
\emph{Amadis de Gaule} appears in the training list, so page-level contamination
is zero; but the \emph{Trésor} excerpts \emph{Amadis}, so on \emph{Amadis}
pages the model has read some of the text it is scored on, in another setting
and another typeface. The report states this and does not correct for it.
